\begin{table}[htbp]
\centering
\caption{Standardized Effect Sizes: UK MOT Mandatory Inspection Threshold}\label{tab:sde}
\begin{threeparttable}
\small
\begin{tabular}{lcccccc}
\toprule
Outcome & $\hat{\beta}$ & SE & SD($Y$) & SDE & SE(SDE) & Class. \\
\midrule
\multicolumn{7}{l}{\textbf{Panel A: Pooled Estimates}} \\
\midrule
First-test failure & 0.0004 & 0.0063 & 0.2849 & 0.001 & 0.022 & Null \\
Second-test failure & 0.0031 & 0.0069 & 0.3023 & 0.010 & 0.023 & Small $+$ \\
\midrule
\multicolumn{7}{l}{\textbf{Panel B: Heterogeneous Effects by Vehicle Fuel Type}} \\
\midrule
First-test failure (petrol) & 0.0024 & 0.0076 & 0.2647 & 0.009 & 0.029 & Small $+$ \\
First-test failure (diesel) & 0.0019 & 0.0129 & 0.3198 & 0.006 & 0.040 & Small $+$ \\
\bottomrule
\end{tabular}
\begin{tablenotes}
\item {\footnotesize \textit{Notes:} \textbf{Country:} United Kingdom. \textbf{Research question:} Does crossing the legally required age-3 MOT vehicle inspection threshold causally reduce defect rates at subsequent annual safety tests for passenger vehicles? \textbf{Policy mechanism:} The UK Ministry of Transport (MOT) test mandates that all passenger vehicles undergo a mandatory roadworthiness inspection at exactly 36 months of age. Vehicles failing the inspection must repair all flagged defects (brakes, lighting, steering, tyres, suspension, emissions) before a re-test pass is issued, creating a legally compelled maintenance intervention at the 36-month birthday. Vehicles may also be tested voluntarily before 36 months without legal requirement. \textbf{Outcome definition:} Binary indicator equal to 1 if the vehicle fails its annual MOT test (any defect requiring repair or advisory note), 0 if the test is passed; measured at first recorded MOT test (Panel A, row 1) or at the second annual test approximately 12 months later (Panel A, row 2, where linked data is available). \textbf{Treatment:} Binary; equal to 1 if the vehicle first tests at or after the 36-month mandatory threshold, 0 if tested voluntarily before 36 months. \textbf{Data:} Driver and Vehicle Standards Agency (DVSA) MOT test results, England, Scotland, and Wales, annual files 2021--2023, accessed via data.gov.uk / S3; sample restricted to vehicles with first test at age 28--46 months (±8-month bandwidth). \textbf{Method:} Sharp regression discontinuity design (RDD) with triangular kernel and MSE-optimal bandwidth (Calonico, Cattaneo, and Titiunik 2014); bias-corrected robust confidence intervals throughout; McCrary density test confirms no manipulation. \textbf{Sample:} All passenger vehicles (Class 4) first registered in Great Britain with a recorded first MOT test, 2022 cohort; multiple registration cohorts for validation; approximately 10\% stratified sample drawn for memory efficiency at 8GB RAM. SDE $= \hat{\beta} / \text{SD}(Y)$ where SD($Y$) is the pre-threshold standard deviation. Classification refers to magnitude, not statistical significance: Large ($|$SDE$| > 0.15$), Moderate (.05$--.15$), Small (.005$--.05$), Null ($< 0.005$).}
\end{tablenotes}
\end{threeparttable}
\end{table}


